\subsection{PRISM}\label{sec:2.2}

% Sources: SYSTEM_PLAN.md section 1; scripts/fig1_prism.py (panel c).
PRISM turns one follow-up scan, the patient's earlier scans and the clinical record into a structured record and a
checked report, in six stages (Figure~\ref{fig:1}c). It takes the four MRI sequences with masks of the tumour, and the
treatment course as of the scan day. From these it computes 27 measurements of the tumour and its change, and the
clinical-history inputs. Three models score the scan, and PRISM states how certain each score is. Deterministic code
then collates every value, with its source, into one record. Finally, a local language model writes the report from
the record, and rule-based checks compare the report with the record. Table~\ref{tab:models} lists PRISM's scores and
the reference scores they are compared with.

% Sources: OUTLINE.md (Fig 1); SYSTEM_PLAN.md section 1.
% FIGURE REDESIGN (2026-11-18), modelled on Tak et al. 2025, NEJM AI, doi:10.1056/AIoa2400703.
% Example paper (local only): docs/literature/examples/tak2025_nejmai.md and tak2025_figures/.
% What we take from it, figure by figure: docs/literature/examples/tak2025_NOTES.md.
% Figure 1a follows their Figure 1A (one real patient on a timeline). Drawn by scripts/fig1_prism.py, which reads the
% MU images from data/mu-glioma-post (patient PatientID_0204; labels by the paper's rule, computed in the script).
\begin{figure}[!htbp]
\centering
\figslot{fig1_system.pdf}{0.6\textheight}{One MU patient's follow-up on real scans; PRISM's segmenter; the method in six stages; the evaluation protocol.}
\caption{\textbf{The task and the method.} (a) One MU-Glioma-Post patient's scans after radiotherapy, placed by weeks
after its end. Each panel is the T1c axial slice with the most enhancement in the supplied mask. The enhancing tumour
is outlined by PRISM's segmenter, a Swin UNETR trained on BraTS 2021 \citep{hatamizadeh2022} (solid), and by the mask
supplied with the dataset (dashed); volumes are the segmenter's, with the supplied mask's in brackets. Enhancement at
3.7 weeks, inside the first 12 weeks, had almost gone by 11.7 weeks without progression; both scans were followed by
at least 120 days without progression and are labelled no progression. The scan at 24.7 weeks was followed by less
than 120 days without progression and is not labelled. The scan at 33.4 weeks, 6 days before the documented
progression date, is marked as progression.\todo{FIGURE: by the paper's rule (first scan on or after the progression
date) this patient's progression scan is the one at 40.9 weeks, and the scan at 33.4 weeks is not labelled; the
script marks a scan within 14 days of the date. Redraw, then reword this sentence.} (b) PRISM's segmenter, shown on
the scan at 33.4 weeks. The Swin UNETR \citep{hatamizadeh2022} takes the four sequences as the channels of one input.
Its Swin Transformer encoder lets each small patch of the image draw on the other patches in its local window
(self-attention) and halves the resolution at each stage; its convolutional decoder restores the resolution, using
skip connections from every stage. Each stage holds two successive Swin blocks (below): layer normalisation (LN),
multi-head self-attention within windows of 7\textsuperscript{3} patches (W-MSA) and a multilayer perceptron (MLP),
each with a residual connection. The second block shifts the windows by half a window (SW-MSA), so that information
passes between neighbouring windows. Three sigmoid outputs (tumour core, whole tumour, enhancing tumour) give the
enhancing, necrosis and FLAIR hyperintensity masks. (c) What PRISM computes at one follow-up scan, in six stages,
drawn for the scans at 24.7 weeks (previous) and 33.4 weeks (current). (1) Inputs: the four sequences, with masks from
PRISM's segmenter or supplied with the dataset (the results use the supplied masks unless stated), and the treatment
course as of the scan day with age, sex, and IDH and MGMT status. (2) Features: 27 measurements of the tumour and its
change ($\mathbf{x}_{\mathrm{img}}$) and the clinical-history inputs ($\mathbf{x}_{\mathrm{hist}}$). (3) Models: three
ridge logistic regressions score the scan from $\mathbf{x}_{\mathrm{img}}$ alone (no dates),
$\mathbf{x}_{\mathrm{hist}}$ alone, and both. (4) Evidence: the outcomes the combined score cannot rule out (no
progression \{0\}, progression \{1\}, or both), chosen by split conformal prediction to include the true outcome for
80\% of scans; whether the imaging and clinical-history scores fall on the same side of their cohort medians; and a
reliability card, which gives the models' results on earlier scans in the same time window and enhancement state. (5)
Collation: deterministic code gathers every score, measurement and statement into one record, each under a fixed key
with its source. (6) Writing and checking: a local language model writes the report from the record alone, each
sentence citing the record keys it uses, and a deterministic checker accepts or rejects the report without repairing
it. Colours mark the models as in every figure. (d) The models were developed on MU-Glioma-Post by five repeats of
fivefold cross validation grouped by patient, fitted once on all 181 labelled MU scans, and applied without retraining
to LUMIERE and UCSF-ALPTDG. The rules were fixed before any score was compared with an external label, except for the
LUMIERE imaging and reduced clinical-history scores, which were added afterwards (Table~\ref{tab:models}).}
\label{fig:1}
\end{figure}

% Drawn by scripts/fig_cohorts.py from the system records and the supplied dates (after Tak et al. 2025, Figure 2).
\begin{figure}[!htbp]
\centering
\figslot{fig_cohorts.pdf}{0.50\textheight}{For MU, LUMIERE and UCSF, the time of each evaluation scan and the interval
since the previous scan, by label.}
\caption{\textbf{When the evaluation scans were taken.} One panel per cohort, headed by its evaluation scans,
patients and share of progression scans. Left, the time of each scan, split by label: at MU-Glioma-Post in weeks
after the end of radiotherapy, with the first 12 weeks shaded; at LUMIERE, which records no radiotherapy dates, in
weeks after surgery; and at UCSF-ALPTDG, which records no end of radiotherapy, in weeks after the first surgery.
Right, weeks since the patient's previous scan. Each column shares one axis across cohorts; scans beyond it are drawn
at the edge and counted. Each point is one scan (filled, progression; open, no progression), and the vertical bar
marks the median. MU shows the 146 scans taken after radiotherapy had started, and UCSF the second scan of each grade 4
patient.}
\label{fig:cohorts}
\end{figure}

% Sources: NUMBERS.md 2.1, 3, 4; DETECTION_TASK.md; UCSF_PLAN.md; CLINICIAN_BASELINE_PLAN.md; CHRISTODOULOU_REPLICATION.md;
% SKELETON_2026-11-20.md (timestamp-only check, scratchpad timing_only.py); METHODS_RESULTS.md E6c (unknown status).
\begin{table}[!htbp]
\centering
\caption[The scores in this paper]{\textbf{The scores in this paper.} Every fitted score except the published
radiomics pipeline is a ridge logistic regression fitted on the training data. Unknown IDH or MGMT status is coded as
wildtype and unmethylated; coding it as its own category moved no AUC by more than 0.04, and every paired CI included
zero. A missing measurement is replaced by the training median, with an indicator that it was missing. Status says
whether a score's own AUC was planned before the analysis or exploratory; the comparisons in Table~\ref{tab:2} were
planned even where a score's AUC on MU was exploratory.}
\label{tab:models}
\footnotesize
\zebrarows
\begin{tabularx}{\linewidth}{@{}R{2.9cm}LR{1.3cm}R{3.4cm}@{}}
\toprule
Score & Inputs & Dates & Status \\
\midrule
Clinical-history model & Age, sex, IDH and MGMT status, days since diagnosis, day of surgery, radiotherapy start, end and dose, fractions, whether chemotherapy had started, the 12-week flag, visit index, biopsy before resection and a previous brain tumour & Yes & Planned \\
Reduced clinical-history model & Age, sex, IDH and MGMT status, days since surgery, a flag for the first weeks after radiotherapy & Partly & Planned at UCSF-ALPTDG, exploratory at LUMIERE \\
Imaging model & 27 measurements of the segmented tumour (volumes, shape and intensity, change since the previous scan, volumes relative to the patient's smallest) & No & Exploratory on MU, planned at UCSF-ALPTDG \\
Images from one scan & The 19 measurements of the current scan & No & Exploratory on MU, planned at UCSF-ALPTDG \\
Combined model & Inputs of both models (externally the reduced clinical-history inputs, volumes and their change, and the interval since a previous scan) & Yes & Exploratory on MU, planned and applied unchanged in its reduced form \\
\addlinespace
Published radiomics pipeline & 256 of 4,892 radiomic features of the current scan ranked within training, class weighted gradient boosting tuned over 400 trials, Platt calibration; the authors' code \citep{christodoulou2026} & No & Planned \\
Time-and-enhancement model & Time window after radiotherapy, log enhancing volume and its change since the previous scan & Yes & Planned \\
Timing alone & Days since the end of radiotherapy or visit number, without a fitted model; or a ridge on these, days since diagnosis and the 12-week flag & Yes & Exploratory \\
Clinical data alone & Age, sex, IDH and MGMT status & No & Planned reference \\
Enhancing volume alone & Log enhancing tumour volume, no fitted model & No & Exploratory reference \\
\bottomrule
\end{tabularx}
\end{table}

% Sources: DETECTION_TASK.md (Plan, B1); UCSF_IMAGING.md; NUMBERS.md 2.1; SEGMENTER_STAGE.md section 2;
% SEGMENTER_BRATS24_PLAN.md; prism/features/reseg.py, reseg_brats24.py, measurements.py; ucsf_imaging.model_inputs;
% refs mahmoud2025, suter2022, fields2024, hatamizadeh2022, deverdier2024, ferreira2024.
\paragraph*{Segmentation and measurements.} PRISM measures the tumour from masks that mark its compartments on every
scan. The masks come either from one of two fixed public segmentation networks, which PRISM runs itself, or from the
masks supplied with each dataset (Table~\ref{tab:1}). The primary results use the masks supplied with each dataset, so
that every volume can be traced to a public file. At MU-Glioma-Post and UCSF-ALPTDG these mark the enhancing tumour,
necrosis, FLAIR hyperintensity and the resection cavity. At LUMIERE they are DeepBraTumIA masks with three labels,
which count the cavity as necrosis.

Both networks are run with their public weights and settings and were never tuned on our cohorts. The first is a Swin
UNETR \citep{hatamizadeh2022} (Figure~\ref{fig:1}b), a network of 62 million parameters that marks the enhancing
tumour, the tumour core and the whole tumour; we use its released model trained on BraTS 2021 (fold 1). It takes
FLAIR, T1c, T1 and T2 as four channels, each scaled to zero mean and unit variance over its non-zero (brain) voxels, and is applied to
blocks of $128^3$ voxels that overlap by half. A voxel belongs to a region when the network's probability exceeds 0.5.
The enhancing tumour is taken directly, necrosis is the core outside the enhancing tumour, and FLAIR hyperintensity is
the rest of the whole tumour. BraTS 2021 contains scans taken before surgery, so this network has no resection-cavity
label and supplies 22 of the 27 measurements. The second network is the winning model of the BraTS 2024 challenge on
adult glioma after treatment \citep{deverdier2024,ferreira2024}, an ensemble of nnU-Net, MedNeXt and Swin UNETR
networks that also marks the resection cavity, and so supplies all 27 measurements. We run it unchanged, with its own
thresholds, from the container distributed by the challenge organisers. LUMIERE's images are stored on an atlas grid whose axes are ordered
differently from the training data of both networks. Before either network, every LUMIERE scan is therefore
reoriented to the order of the training data, without resampling, and the label map is turned back. How each network
changes the results is tested in Section~\ref{sec:2.3}.

From the masks and the T1c and FLAIR images, PRISM computes 27 measurements (Supplementary Table~S7). Nineteen describe
the current scan: the log volumes of the enhancing tumour, necrosis, FLAIR hyperintensity, resection cavity, tumour
core and whole lesion; eight measures of composition, shape and position relative to the cavity; and five measures of
T1c and FLAIR intensity relative to normal-appearing brain, taken as the brain more than 10~mm from the lesion and the
cavity. Eight describe change: the log ratio of each of the six volumes to the previous scan, and of the enhancing
tumour and whole lesion to the patient's smallest volume so far. Change uses earlier scans only and is set to zero at
a patient's first scan. Five measurements need the cavity: its volume, whether there is one, the share of enhancing
tumour within 5~mm of it, the distance of the enhancing tumour from it, and its change since the previous scan.

The clinical-history inputs describe the treatment course as of the scan day: days since diagnosis and surgery,
radiotherapy start, end and dose, and whether chemotherapy had started. PRISM also places the scan in a time window
after radiotherapy (before its end, within 12 weeks after its end, later, or unknown when dates are missing), and it
takes age, sex, and IDH mutation and MGMT promoter methylation status. Each input enters only once it is known at the
scan: the prescribed dose from the start of radiotherapy, its duration and the time since its end only after it has
ended, and IDH and MGMT status only from the diagnosis date (they count as unknown at 8 grade 4 scans taken before
it, none of them in the evaluation set). The three models in the next paragraph use these inputs and the 27 measurements.

% Sources: prism/models/formulations.py; prism/experiments/c3_radiomics.py (BASE_PRE); NUMBERS.md 2.1;
% DETECTION_TASK.md (External test); SYSTEM_RESULTS.md (Deviations item 1).
